\subsection*{Stability and reproducibility of GAN training}
The instability of adversarial training is a long-standing concern.
Heusel et al.~\cite{heusel2017fid} proposed the Fr\'echet Inception
Distance and a two-timescale update rule.
Borji~\cite{borji2022pros} surveys the proliferation of evaluation
metrics that followed. None of these directly target seed-to-seed
reproducibility of training behavior on a \emph{fixed} dataset; they
are distributional similarity measures evaluated on a single trained
model. Lucic et al.~\cite{lucic2018gans} argued that hyperparameter
sensitivity dominates apparent GAN improvements and multi-seed
training is the only honest way to compare architectures; their
analysis aggregates seeds within a single dataset, while we explicitly
compute per-dataset seed-variance and use it as a discriminative
scalar \emph{across} datasets. The four-axis Stability Score we
describe in Section~\ref{sec:methods} is, to our knowledge, the first
such aggregate in the financial time series setting.

\subsection*{Microstructural taxonomy}
Briola et al.~\cite{briola2024} proposed a static taxonomy of equity
LOBs into large-/small-tick and high-/low-information regimes based on
spread, queue imbalance and price-impact statistics computed across
three years of NASDAQ data. The taxonomy presumes the macroscopic
stability of spread distributions over a multi-year window. We argue
that this assumption fails in crypto: spreads transition from
large-tick to small-tick behavior within seconds during liquidation
cascades, so any taxonomy that aspires to clock-invariance must be
dynamic, regime-aware, and computed on a rolling basis. The
Stationarity Score we introduce is one such dynamic analogue and is,
in fact, a crypto-adapted refactoring of Briola et al.'s intent.
